% Generated draft: requires mathematical and source review.
\documentclass[11pt]{article}
\usepackage[margin=25mm]{geometry}
\usepackage{fontspec}
\setmainfont{Latin Modern Roman}
\usepackage{amsmath,amssymb,mathtools,unicode-math}
\setmathfont{Latin Modern Math}
\usepackage{enumitem,needspace,xurl,hyperref}
\setlength{\parindent}{0pt}
\setlength{\parskip}{5pt}
\setlength{\emergencystretch}{3em}
\newcommand{\N}{\mathbb{N}}
\newcommand{\Z}{\mathbb{Z}}
\newcommand{\Q}{\mathbb{Q}}
\newcommand{\R}{\mathbb{R}}
\newcommand{\C}{\mathbb{C}}
\newcommand{\F}{\mathbb{F}}
\newcommand{\A}{\mathbb{A}}
\newcommand{\PP}{\mathbb{P}}
\newcommand{\E}{\mathbb{E}}
\newcommand{\eps}{\varepsilon}
\newcommand{\dd}{\,\mathrm{d}}
\begin{document}
\section*{UnsolvedMath ID: 11000065}
\textbf{Source problem code:} catalogue entry
\subsection{Definitions and mathematical statement}
\textbf{Problem 5.21 (Malcev mod $2$).} Give an explicit finite presentation for the $\F_2$-Lie algebra denoted $L_2(\mathcal I_{g,1})$ in the source. Here $\mathcal I_{g,1}$ is the corresponding Torelli group, but the extracted statement does not specify whether the subscript $1$ denotes a marked point or a boundary component, nor does it define the functor $L_2$ or give a genus range.

For an ordinary Lie algebra, an explicit finite presentation would consist of finite sets of generators $X$ and Lie relations $R$ together with an identified isomorphism
\[
 L_2(\mathcal I_{g,1})\cong\operatorname{Lie}_{\F_2}(X)/\langle R\rangle_{\mathrm{Lie}}.
\]
If the intended algebra is restricted, graded, or completed, the associated structure, grading, or completion and the appropriate notion of presentation must also be specified. In characteristic two, the convention that an ordinary Lie bracket is alternating is relevant. The record does not decide whether $L_2$ uses an ordinary lower central series tensored with $\F_2$, a mod-$2$ central series, or a pro-$2$/restricted construction; these cannot be interchanged silently.

The neighboring source record 11000064 stresses that reduction of the characteristic-zero information captures only part of $L_2(\mathcal I_{g,1})$. It singles out additional $2$-torsion, written $B_2$, associated with the Rochlin invariant rather than merely the algebraic Johnson-homomorphism contribution. Thus the requested presentation is not just an unsupported reduction modulo $2$ of a characteristic-zero presentation.

The appended paragraph introduces a separate real Lie algebra from the Johnson filtration:
\[
 \mathfrak h_g=\bigoplus_{k=1}^{\infty}
 \left(\mathcal I_g(k)/\mathcal I_g(k+1)\right)\otimes_{\Z}\R,
\]
with bracket induced by group commutators. It reports that Hain, in section 14 of [Ha3], proves that the Johnson filtration is not cofinal with the lower central series of $\mathcal I_g$ and relates $\mathfrak h_g$ to the Malcev Lie algebra $\mathfrak t_g$. This paragraph is context leading to a following question, not a replacement of the mod-$2$ presentation problem by a real-coefficient one.
\subsection{Short English statement}
Find finite generators and relations for the intended characteristic-two Lie algebra of the marked or bordered Torelli group, retaining its extra $2$-torsion information. The definition of that Lie algebra and the surface convention are missing.
\subsection{Sources}
\begin{itemize}
\item Local UnsolvedMath registry, permanent ID 11000065, source code catalogue entry; no individual proposer or proposal year supplied.
\item Farb, edited volume \emph{Problems on Mapping Class Groups and Related Topics}, Problem 5.21, PDF page 51; \url{https://www.math.uchicago.edu/~farb/papers/mcgbook.pdf}. The registry records PDF-font-structure extraction and access on 2026-07-29; its inventory distinguishes 329 labeled headings by source page.
\item Adjacent local ID 11000064 supplies the discussion of characteristic-zero reduction, $B_2$, and the Rochlin invariant. Source label [Ha3], section 14, is named in the present excerpt, but its complete bibliography entry is absent.
\end{itemize}
\end{document}
